\documentclass[10pt,a4paper,sans]{moderncv}
\moderncvstyle{classic}
\moderncvcolor{blue}
\usepackage[utf8]{inputenc}
\usepackage[scale=0.85]{geometry}
\usepackage{helvet}

\name{Alexandro}{Disla}
\title{Responsable MEAL --- Suivi, Évaluation, Redevabilité et Apprentissage}
\address{Rue Saint-Louis Jeanty \#14, Route de Frère}{Port-au-Prince, Haïti}{}
\phone[mobile]{(+509) 4148-3700}
\email{alexandrodisla@hotmail.com}
\extrainfo{GitHub: \url{https://github.com/AD0791}}

\begin{document}
\makecvtitle

\section{Profil Professionnel}
\cvitem{}{Plus de dix ans en suivi, évaluation et systèmes d'information --- dont \textbf{cinq ans de MEAL en ONG} sur trois secteurs en Haïti (santé et VIH, eau et assainissement, éducation) et \textbf{huit ans de suivi de l'exécution et de modélisation statistique} à la Direction de Planification Économique et Sociale du MPCE. Mon travail porte sur ce qui rend les chiffres d'un programme défendables : cadres logiques et tableaux de suivi des indicateurs, protocoles d'assurance qualité des données établis puis supervisés, collecte numérique standardisée pour que la donnée soit juste à la saisie plutôt que réparée ensuite. Économiste appliqué de formation et ingénieur de données en exercice, je construis et administre moi-même les bases, les pipelines et les tableaux de bord dont dépend un bureau pays --- je ne me limite pas à les spécifier. Résidant à Port-au-Prince, disponible pour les déplacements terrain ; français, anglais et créole haïtien courants.}

\section{Compétences Clés}
\cvitem{}{
\begin{itemize}
    \item \textbf{Pilotage d'un cadre MEAL \& harmonisation :} Stratégies et cadres de suivi-évaluation, chaînes de résultats, cadres logiques, tableaux de suivi des indicateurs (ITT), plans MEAL, valeurs de référence, cibles et désagrégations exigées --- dont la désagrégation par sexe et par âge posée dès la définition de l'indicateur. Standardisation des formulaires et des flux de données entre sites et entre secteurs.
    \item \textbf{Qualité des données \& contrôle :} Protocoles DQA établis et supervisés, vérification par remontée aux registres sources, suivi des actions correctives jusqu'à clôture ; dédoublonnage des listes de bénéficiaires et fiabilisation des décomptes servant au reporting bailleur et aux propositions.
    \item \textbf{Systèmes d'information \& reporting bailleurs :} Bases relationnelles et pipelines ETL conçus, administrés et sécurisés (contrôle d'accès par rôle, stockage cloud) ; tableaux de bord de suivi hebdomadaire ; production et soumission d'indicateurs PEPFAR/MER au bailleur USAID via les plateformes DHIS2/DATIM \textit{(en tant que producteur de données)} ; projet sous contrat des Nations Unies (UNOPS).
    \item \textbf{Encadrement, apprentissage \& coordination :} Supervision d'agents de collecte sur programmes multi-sites, encadrement technique de personnel de suivi-évaluation, conception et animation de formations aux outils et à la qualité des données ; coordination avec les partenaires gouvernementaux (MSPP, DINEPA, MEF) et restitution des constats techniques à une direction non technique.
\end{itemize}}

\section{Expériences MEAL et Qualité de Programme}

\cventry{Jan 2026 -- Mars 2026}{Analyste de Données \& Consultant MEAL}{Anseye Pou Ayiti}{Haïti}{}{
\begin{itemize}
    \item Définition de la stratégie de collecte et mise en œuvre des indicateurs de performance de trois programmes de leadership éducatif, de la définition de l'indicateur jusqu'à sa restitution.
    \item Codage des instruments XLSForm pour ODK et Google Forms avec logiques de saut et contraintes de validation, afin que les erreurs soient arrêtées à la saisie et non corrigées en aval.
    \item Administration du système ODK et du panneau d'administration, maintenance des scripts d'intégration Python sur AWS, et formation des agents de terrain aux outils de collecte.
\end{itemize}}

\cventry{Oct 2023 -- Jan 2026}{Consultant mWater \& Analyste de Données}{HANWASH}{Haïti / Télétravail}{}{
\begin{itemize}
    \item \textbf{Conception MEAL :} Raffinement des plans MEAL et des tableaux de suivi des indicateurs alignés sur les normes JMP (OMS/UNICEF) pour des programmes d'eau, d'hygiène et d'assainissement.
    \item \textbf{Gestion des données :} Déploiement d'enquêtes géospatiales complexes sur mWater, administration de la base et résolution des incidents de collecte terrain pour préserver l'intégrité des séries.
    \item \textbf{Automatisation \& capacités :} Tableaux de bord automatisés (mWater, Power BI) connectés aux API de collecte pour un reporting en temps réel ; rédaction des guides techniques et formation des équipes à leur usage autonome.
\end{itemize}}

\cventry{Oct 2021 -- Août 2024}{Officier de Suivi \& Évaluation}{Impact Youth Project, Caris Foundation International}{Haïti}{}{
\begin{itemize}
    \item \textbf{Systèmes de collecte :} Mise en place de la collecte numérique (CommCare, HIVHaiti) sur un programme de santé et de VIH multi-sites, avec standardisation des formulaires et des flux entre sites.
    \item \textbf{Assurance qualité (DQA) :} Établissement et supervision des protocoles de Data Quality Assessment sur les données intégrées MySQL --- remontée des chiffres rapportés jusqu'aux registres sources, et suivi des actions correctives jusqu'à clôture plutôt que jusqu'à leur simple consignation.
    \item \textbf{Reporting bailleur :} Production et soumission des indicateurs PEPFAR/MER destinés à l'USAID via DHIS2/DATIM, avec les contrôles de cohérence et le respect du cycle de rapportage que cela impose.
    \item \textbf{Encadrement :} Supervision des agents de collecte sur l'ensemble des sites et encadrement technique du personnel de suivi-évaluation sur les procédures de collecte et les contrôles qualité.
    \item \textbf{Automatisation :} Rapports d'activité développés en Python et SQL, réduisant le temps de traitement manuel de 40 \% ; tableaux de bord programme maintenus sur AWS (Quarto) et Power BI pour le suivi hebdomadaire des indicateurs.
\end{itemize}}

\cventry{Nov 2015 -- Août 2023}{Économiste --- Direction de Planification Économique et Sociale (DPES)}{Ministère de la Planification et de la Coopération Externe (MPCE)}{Port-au-Prince}{}{
\begin{itemize}
    \item Contribution aux rapports de suivi de l'exécution du Plan Triennal d'Investissement 2014--2016 : suivi d'indicateurs de projets d'investissement public sur l'ensemble des secteurs et des départements géographiques.
    \item Affectation à l'unité statistique et modélisation, et participation à la construction du << Modèle de Planification du Développement >> --- outil de prévision, de simulation et d'analyse d'impact des politiques publiques.
    \item Participation aux réunions de cadrage macroéconomique avec le Ministère de l'Économie et des Finances et suivi des indicateurs macroéconomiques nationaux.
\end{itemize}}

\section{Expériences en Ingénierie de Données et Systèmes}

\cventry{Févr 2026 -- Présent}{Ingénieur Logiciel (Fullstack)}{Tekkod LLC}{Télétravail}{}{Direction d'un programme de modernisation d'infrastructure mobile avec fiabilité hors ligne et sécurité des données pour un usage terrain ; tableaux de bord Looker Studio avec recherche et filtrage avancés pour la visibilité opérationnelle.}

\cventry{Mars 2021 -- Mai 2024}{Développeur Backend \& Ingénieur de Données}{Tekkod LLC}{Télétravail}{}{Pipelines ETL (Apache Beam) migrant les données MongoDB vers BigQuery et administration de bases relationnelles en accès concurrent ; services REST sécurisés (FastAPI, JWT) avec contrôle d'accès par rôle et documentation OpenAPI complète.}

\cventry{Nov 2015 -- Mars 2021}{Consultant Logiciel \& Spécialiste Intégration de Données}{CassionSoft (Projet UNOPS)}{Haïti}{}{Intégration de données sur un projet sous contrat des Nations Unies, avec exposition directe aux procédures et aux standards documentaires onusiens ; authentification RBAC (OAuth2/JWT) et validation de la logique d'intégration par tests unitaires.}

\section{Formation}
\cvitem{2010--2014}{Études Supérieures en Économie Appliquée --- C.T.P.E.A (Centre de Techniques de Planification et d'Économie Appliquée), Port-au-Prince. \textit{Scolarité complétée ; Diplôme d'Études Supérieures en attente de la soutenance du mémoire de sortie --- attestation officielle jointe au dossier.}}
\cvitem{2009--2010}{Baccalauréat (Mention) --- Institution Saint-Louis de Gonzague, Port-au-Prince}
\cvitem{Autoformation}{FlatIron School (Software Engineering), CodeWithMosh, Udemy}

\section{Compétences Techniques}
\cvitem{Outils MEAL}{\textbf{KoboToolbox, ODK, XLSForm}, CommCare, mWater, HIVHaiti, DHIS2/DATIM \textit{(soumission de données)}, cadres logiques, ITT, protocoles DQA}
\cvitem{Analyse}{\textbf{Excel avancé}, \textbf{Power BI (expert)}, SQL, Python (Pandas), R (RMarkdown/Quarto), Looker Studio, modélisation statistique, échantillonnage et analyse de tendances}
\cvitem{Infrastructure}{MySQL, PostgreSQL, MongoDB, Azure, AWS, Google Cloud (BigQuery), Docker, Git}
\cvitem{Langues}{\textbf{Français (maternel)}, \textbf{Anglais (courant)}, \textbf{Créole haïtien (maternel)}}

\end{document}
